% !TeX root = ../机器学习教材（李航）笔记.tex
% 由主文档通过 \input 载入；请编译主文档。
\setcounter{chapter}{18} % 保留原书章号。
\chapter{概率潜在语义分析}
\label{ch:lh-plsa}

一个文本可以同时谈论多个话题，同一个单词也可能出现在不同话题中。因此，只把文本或单词划分到互不相交的类别，往往不足以表达它们的语义联系。\emph{概率潜在语义分析}（probabilistic latent semantic analysis，PLSA），又称\emph{概率潜在语义索引}（probabilistic latent semantic indexing，PLSI），用不可直接观测的话题连接文本与单词：每个文本对应一个话题分布，每个话题对应一个单词分布，从可观测的单词频数反推这两类分布。这是一种无监督学习方法，话题数由建模者预先指定，话题的语义名称则需要根据学到的分布解释，不能把名称本身当作训练标签。

本章先建立生成模型及其等价的共现模型，再从参数、几何与矩阵分解三个角度解释模型，最后推导极大似然估计的期望极大算法（expectation-maximization algorithm，EM）。所需的条件概率、有限混合分布、拉格朗日乘子与对数函数凹性都将在具体推导中使用；全文采用自然对数，明确区分观测频数、模型概率与隐变量的后验概率。

\section{概率潜在语义分析模型}
\label{sec:lh-plsa-model}

\subsection{基本想法}
\label{subsec:lh-plsa-idea}

设词表为 $\mathcal W=\{w_1,\ldots,w_M\}$，文本指标集合为 $\mathcal D=\{d_1,\ldots,d_N\}$，预先给定的话题集合为 $\mathcal Z=\{z_1,\ldots,z_K\}$。这里 $M,N,K$ 都是正整数，$M$ 计算不同词的种数，而非语料中单词出现的总次数。用 $n_{ij}=n(w_i,d_j)\in\{0,1,2,\ldots\}$ 表示单词 $w_i$ 在文本 $d_j$ 中出现的次数，得到\emph{单词--文本矩阵}
\begin{equation}
X=(n_{ij})\in\mathbb R^{M\times N},\qquad
L_j=\sum_{i=1}^M n_{ij},\qquad
n=\sum_{j=1}^N L_j>0.
\label{eq:lh-plsa-counts}
\end{equation}
$L_j$ 为文本长度，$n$ 为语料中词项出现的总次数。以下删去空文本和全零词行，因而 $L_j>0$，且词表中的每个词至少出现一次；文本可以有不同长度。

模型只使用一个词出现在哪个文本以及出现多少次，不使用词序，这称为\emph{词袋假设}（bag-of-words assumption）。话题并非预先列出的关键词集合，而是整个词表上的概率分布；一个词可以在多个话题下具有正概率。相应地，文本也不必只属于一个话题，而可以具有不同的话题混合比例。这种表示实现了对文本和单词的软聚类（soft clustering），但“同属一个话题”表示模型给出的统计关联，不直接等同于词义相同。

为了区分随机变量与其取值，以下用 $D$、$W$、$Z$ 分别表示文本指标、单词和话题随机变量，其取值分别为 $d_j$、$w_i$、$z_k$。一次观测对应一个词项出现所给出的二元组 $(w_i,d_j)$，与之相关的话题 $z_k$ 不可直接观测。\emph[caution]{隐变量是每次词项出现所对应的话题，而不是整个文本唯一的话题。}同一文本内不同词项出现可以由不同话题产生。

\subsection{生成模型}
\label{subsec:lh-plsa-generative}

\begin{definition}[概率潜在语义分析的生成模型]
\label{def:lh-plsa-generative}
设
\[
\pi_j=P(D=d_j),\qquad
\theta_{kj}=P(Z=z_k\mid D=d_j),\qquad
\phi_{ik}=P(W=w_i\mid Z=z_k).
\]
记 $\pi=(\pi_1,\ldots,\pi_N)^{\mathsf T}$。各参数非负，并满足
\begin{equation}
\sum_{j=1}^N\pi_j=1,\qquad
\sum_{k=1}^K\theta_{kj}=1\quad(1\leq j\leq N),\qquad
\sum_{i=1}^M\phi_{ik}=1\quad(1\leq k\leq K).
\label{eq:lh-plsa-normalization}
\end{equation}
假设给定话题后单词与文本条件独立，即 $W\perp D\mid Z$，并规定
\begin{equation}
P(D=d_j,Z=z_k,W=w_i)=\pi_j\theta_{kj}\phi_{ik}.
\label{eq:lh-plsa-complete-joint}
\end{equation}
将话题边缘化，得到生成模型的观测概率
\begin{equation}
\begin{aligned}
p_{ij}&=P(W=w_i,D=d_j)
=\pi_j\sum_{k=1}^K\phi_{ik}\theta_{kj},\\
p_{i\mid j}&=P(W=w_i\mid D=d_j)
=\sum_{k=1}^K\phi_{ik}\theta_{kj},
\end{aligned}
\label{eq:lh-plsa-generative}
\end{equation}
其中条件概率式用于 $\pi_j>0$ 的文本。
\end{definition}

对于一次词项出现，生成过程依次为：按 $\pi$ 选取文本指标 $d_j$，按该文本的话题分布 $\theta_{\cdot j}$ 选取话题 $z_k$，再按该话题的词分布 $\phi_{\cdot k}$ 选取单词 $w_i$；观测时只保留文本与单词。图模型的依赖关系可写成 $D\to Z\to W$。条件独立假设具体意味着
\[
P(W=w_i\mid Z=z_k,D=d_j)=\phi_{ik},
\]
即给定话题后，同一个话题的词分布在所有文本中共享，而不是每个文本另有一套话题词分布。

也可以固定每个文本的长度 $L_j$，对该文本的每个词项出现独立抽取话题和单词。这时该文本的词频向量服从多项分布，单次抽取所用的是词表上的类别分布（categorical distribution）。将单次类别抽取与多次抽取得到的多项计数区分开，有助于理解概率参数和观测频数的关系。两种观察方式分别给出联合似然和给定文本指标后的条件似然；在忽略与参数无关的多项系数后，它们为
\begin{equation}
\mathcal L_{\mathrm{joint}}=\prod_{i=1}^M\prod_{j=1}^N
\left(\pi_j p_{i\mid j}\right)^{n_{ij}},\qquad
\mathcal L_{\mathrm{cond}}=\prod_{i=1}^M\prod_{j=1}^N
p_{i\mid j}^{n_{ij}}.
\label{eq:lh-plsa-likelihoods}
\end{equation}
若把有顺序的独立词项样本汇总为计数，还要乘以相应的多项系数；该系数不影响极大似然估计。这里的乘积分解依赖词项独立抽样的建模假设，不能仅由条件独立 $W\perp D\mid Z$ 单独推出。

\subsection{共现模型}
\label{subsec:lh-plsa-cooccurrence}

生成模型从文本出发选择话题，共现模型则先选择话题，再分别生成文本和单词。在后一种表示下，文本与单词的地位是对称的，话题解释了两者的共现关联。

\begin{definition}[概率潜在语义分析的共现模型]
\label{def:lh-plsa-cooccurrence}
设 $\rho_k=P(Z=z_k)$、$\psi_{jk}=P(D=d_j\mid Z=z_k)$，并保留 $\phi_{ik}=P(W=w_i\mid Z=z_k)$，记 $\rho=(\rho_1,\ldots,\rho_K)^{\mathsf T}$。各参数非负，满足
\[
\sum_{k=1}^K\rho_k=1,\qquad
\sum_{j=1}^N\psi_{jk}=1,\qquad
\sum_{i=1}^M\phi_{ik}=1.
\]
在 $W\perp D\mid Z$ 的假设下，共现概率为
\begin{equation}
p_{ij}=\sum_{k=1}^K\rho_k\phi_{ik}\psi_{jk}.
\label{eq:lh-plsa-cooccurrence}
\end{equation}
其生成过程为先按 $\rho$ 选取话题，再在给定话题后独立地按 $\psi_{\cdot k}$、$\phi_{\cdot k}$ 选取文本与单词，对应依赖关系 $D\leftarrow Z\to W$。
\end{definition}

\begin{proposition}[两种模型的等价性]
\label{prop:lh-plsa-equivalence}
生成模型与共现模型描述相同的一族单词--文本联合分布。对于正概率的条件事件，两者可由贝叶斯公式相互转换；零概率条件事件上的条件分布可以任意指定，不改变联合分布。
\end{proposition}
\begin{proof}
给定生成模型的参数，令
\begin{equation}
\rho_k=\sum_{j=1}^N\pi_j\theta_{kj},\qquad
\psi_{jk}=\frac{\pi_j\theta_{kj}}{\rho_k}\quad(\rho_k>0).
\label{eq:lh-plsa-generative-to-cooccurrence}
\end{equation}
由非负性与归一化条件，$\sum_k\rho_k=1$，且每个正权重话题满足 $\sum_j\psi_{jk}=1$。若 $\rho_k=0$，则所有 $\pi_j\theta_{kj}=0$，可任取一个归一化的 $\psi_{\cdot k}$。于是总有 $\rho_k\psi_{jk}=\pi_j\theta_{kj}$，代入式~\eqref{eq:lh-plsa-cooccurrence} 即得式~\eqref{eq:lh-plsa-generative}。

反过来，给定共现模型，令
\begin{equation}
\pi_j=\sum_{k=1}^K\rho_k\psi_{jk},\qquad
\theta_{kj}=\frac{\rho_k\psi_{jk}}{\pi_j}\quad(\pi_j>0).
\label{eq:lh-plsa-cooccurrence-to-generative}
\end{equation}
同理可验证归一化条件；若 $\pi_j=0$，任取归一化的 $\theta_{\cdot j}$，仍有 $\pi_j\theta_{kj}=\rho_k\psi_{jk}$。因此两种表示给出相同的联合分布。\qedhere
\end{proof}

这也完成了教材中证明两种模型等价的习题。等价的是可表示的概率分布，而不是各个参数数组逐项相同：$\theta_{kj}$ 描述给定文本后的话题概率，$\psi_{jk}$ 描述给定话题后的文本概率，不能交换两者的条件方向。生成模型便于解释文本的混合组成，共现模型便于解释文本与单词如何通过话题形成关联。

\subsection{模型性质}
\label{subsec:lh-plsa-properties}

直接为 $MN$ 个单词--文本对指定联合概率，需要 $MN-1$ 个自由参数。生成模型的参数化在概率单纯形内部需要
\begin{equation}
(N-1)+N(K-1)+K(M-1)=K(M+N-1)-1
\label{eq:lh-plsa-parameter-count}
\end{equation}
个自由坐标；共现模型相应为 $(K-1)+K(M-1)+K(N-1)$，结果相同。若文本边缘概率 $\pi$ 已固定，则生成模型只需后两组参数。当 $K$ 远小于 $M$ 和 $N$ 时，参数规模由 $O(MN)$ 降为 $O(K(M+N))$。这是参数化的坐标数，不表示不同参数一定对应不同的联合分布；例如，统一置换话题指标不会改变模型。较少的参数限制了模型容量，但不能据此断言不会过拟合。

几何上，用\emph{概率单纯形}（probability simplex）
\begin{equation}
\Delta^{M-1}=\left\{v\in\mathbb R^M:
v_i\geq0,\ \sum_{i=1}^M v_i=1\right\}
\label{eq:lh-plsa-simplex}
\end{equation}
表示词表上的全部概率分布。每个话题对应点 $\phi_{\cdot k}\in\Delta^{M-1}$，每个文本的模型词分布满足
\begin{equation}
p_{\cdot\mid j}=\sum_{k=1}^K\theta_{kj}\phi_{\cdot k}
\in\operatorname{conv}\{\phi_{\cdot1},\ldots,\phi_{\cdot K}\},
\label{eq:lh-plsa-convex-hull}
\end{equation}
其中 $\operatorname{conv}$ 表示凸包（convex hull）。非负且和为 $1$ 的系数使这里的线性组合成为凸组合，所以所有文本的模型分布都被限制在话题点的凸包中；观测相对频率 $n_{\cdot j}/L_j$ 则未必恰好落在其中，学习是在寻找合适的近似。$K=1$ 时凸包为一点；$K=2$ 且两个话题不同，凸包为线段。一般地，凸包维数至多为 $\min\{K-1,M-1\}$，只有 $K$ 个话题点仿射无关，即差向量 $\phi_{\cdot2}-\phi_{\cdot1},\ldots,\phi_{\cdot K}-\phi_{\cdot1}$ 线性无关时，才是 $(K-1)$ 维单纯形。此时混合系数是唯一的重心坐标；若存在仿射相关性，同一个词分布可能有不同的话题混合表示。教材所称的“话题单纯形”应结合这个条件理解。

矩阵形式使模型与潜在语义分析的联系更清楚。令
\[
\Phi=(\phi_{ik})\in\mathbb R^{M\times K},\quad
\Theta=(\theta_{kj})\in\mathbb R^{K\times N},\quad
\Psi=(\psi_{jk})\in\mathbb R^{N\times K},
\]
并令 $B=(p_{i\mid j})$、$A=(p_{ij})$，则
\begin{equation}
B=\Phi\Theta,\qquad
A=\Phi\Theta\operatorname{diag}(\pi)
=\Phi\operatorname{diag}(\rho)\Psi^{\mathsf T}.
\label{eq:lh-plsa-matrix-factorization}
\end{equation}
$\operatorname{diag}(v)$ 表示以向量 $v$ 为对角元的对角矩阵。$\Phi$、$\Theta$、$\Psi$ 都是非负且列归一化的矩阵，$\operatorname{diag}(\rho)$ 的非负对角元之和为 $1$，因此 $B$ 的各列是条件分布，$A$ 的全部元素之和为 $1$。二者的秩都不超过 $K$，但它们分别是条件概率矩阵与联合概率矩阵，不能直接与未归一化的频数矩阵 $X$ 混同。

\emph{潜在语义分析}（latent semantic analysis，LSA）通常对单词--文本矩阵作截断奇异值分解：取 $K\leq\min\{M,N\}$，保留最大的 $K$ 个奇异值及其左右奇异向量，得到 $X\approx U_K\Sigma_KV_K^{\mathsf T}$。其中左右奇异向量分别正交，但其分量可以为负。PLSA 则在非负、归一化约束下估计概率因子，不要求正交，学习目标来自似然。因此，两者虽然都通过低维因子描述文本与单词的关联，却不能把奇异向量直接当作话题概率，也不能把截断奇异值分解直接当作 PLSA 的参数估计。

\section{概率潜在语义分析的算法}
\label{sec:lh-plsa-algorithm}

\subsection{极大似然目标与隐变量}
\label{subsec:lh-plsa-objective}

根据式~\eqref{eq:lh-plsa-likelihoods}，联合对数似然在略去数据常数后分解为
\begin{equation}
\begin{aligned}
\ell_{\mathrm{joint}}(\pi,\Phi,\Theta)
&=\sum_{j=1}^N L_j\log\pi_j+\ell(\Phi,\Theta),\\
\ell(\Phi,\Theta)&=\sum_{i=1}^M\sum_{j=1}^N
n_{ij}\log\left(\sum_{k=1}^K\phi_{ik}\theta_{kj}\right).
\end{aligned}
\label{eq:lh-plsa-loglikelihood}
\end{equation}
以下对零频数项约定 $0\log p=0$；若 $n_{ij}>0$ 而模型概率为 $0$，则对数似然为 $-\infty$。学习是在式~\eqref{eq:lh-plsa-normalization} 的非负、归一化约束下极大化该目标。对 $\pi$ 的优化与其他两组参数分离，拉格朗日方程为 $L_j/\pi_j=\lambda$，求和得 $\lambda=n$，从而
\begin{equation}
\widehat\pi_j=\frac{L_j}{n}.
\label{eq:lh-plsa-document-marginal}
\end{equation}
由于各 $L_j>0$，目标 $\sum_jL_j\log\pi_j$ 在单纯形内部严格凹，并在边界趋于 $-\infty$，该解是唯一最大点。因此 $\widehat\pi_j$ 是随机词项出现来自文本 $d_j$ 的经验概率，只有文本等长时才为 $1/N$；后续只需优化 $\Phi$ 与 $\Theta$。

似然准则也可以从分布拟合理解。定义经验联合分布 $\widehat p_{ij}=n_{ij}/n$，则经验分布到模型分布的\emph{库尔贝克--莱布勒散度}（Kullback--Leibler divergence，KL divergence）为
\begin{equation}
D_{\mathrm{KL}}(\widehat p\|p)
=\sum_{i,j:\,\widehat p_{ij}>0}\widehat p_{ij}
\log\frac{\widehat p_{ij}}{p_{ij}}.
\label{eq:lh-plsa-kl}
\end{equation}
当经验正概率处的模型概率均为正时，由 $-\log x\geq1-x$ 可得该散度非负；若某处模型概率为零，则散度为 $+\infty$。由于
\[
\frac1n\ell_{\mathrm{joint}}
=\sum_{i,j}\widehat p_{ij}\log\widehat p_{ij}
-D_{\mathrm{KL}}(\widehat p\|p),
\]
极大似然估计等价于在 PLSA 可表示的分布族中最小化该散度。这个等价关系说明了概率拟合与平方误差矩阵近似的区别，而不提供整个问题的凸性或全局解保证。

直接对 $\ell$ 求极值会遇到“对数内含求和”的耦合。若知道每次词项出现来自哪个话题，令 $n_{ijk}$ 表示单词--文本对 $(w_i,d_j)$ 中来自话题 $z_k$ 的次数，则 $\sum_kn_{ijk}=n_{ij}$，完全数据的参数相关对数似然为
\begin{equation}
\ell_c=\sum_jL_j\log\pi_j+
\sum_{i,j,k}n_{ijk}(\log\phi_{ik}+\log\theta_{kj}).
\label{eq:lh-plsa-complete-loglikelihood}
\end{equation}
这时不同概率列的估计可以分开完成。EM 的思路是在当前参数下求出未知计数的条件期望，再以这些期望计数重新估计参数。隐藏的数据是话题取值或对应的计数 $n_{ijk}$，后验概率只是计算条件期望的工具，不能把一个概率分布本身说成观测数据或隐数据。

\subsection{E 步与 M 步的推导}
\label{subsec:lh-plsa-em-derivation}

记 $\eta=(\Phi,\Theta)$，第 $t$ 轮参数为 $\eta^{(t)}$。\emph{E 步}使用当前参数计算\emph{责任度}（responsibility），即每个已观测词项出现来自话题 $z_k$ 的后验概率：
\begin{equation}
r_{ijk}^{(t)}=P_{\eta^{(t)}}(Z=z_k\mid W=w_i,D=d_j)
=\frac{\phi_{ik}^{(t)}\theta_{kj}^{(t)}}
{\sum_{h=1}^K\phi_{ih}^{(t)}\theta_{hj}^{(t)}}.
\label{eq:lh-plsa-e-step}
\end{equation}
其中 $\pi_j$ 在分子分母中约去，责任度同时依赖单词和文本，通常并不等于文本先验话题比例 $\theta_{kj}$。只需对 $n_{ij}>0$ 的位置计算它；这些位置的分母必须为正。由给定观测后的条件分布，
\begin{equation}
\mathbb E_{\eta^{(t)}}[n_{ijk}\mid X]=n_{ij}r_{ijk}^{(t)}.
\label{eq:lh-plsa-expected-count}
\end{equation}
同一个单词--文本对可能出现多次，故这里必须乘以频数 $n_{ij}$，而不能把每个非零矩阵元都只计为一次样本。

将式~\eqref{eq:lh-plsa-complete-loglikelihood} 对隐数据取上述条件期望，并去掉与 $\eta$ 无关的项，得到
\begin{equation}
Q(\eta\mid\eta^{(t)})=
\sum_{i,j,k}n_{ij}r_{ijk}^{(t)}
(\log\phi_{ik}+\log\theta_{kj}).
\label{eq:lh-plsa-q-function}
\end{equation}
\emph{M 步}在非负和列归一化约束下极大化 $Q$。注意 $r^{(t)}$ 在此步骤中已经固定，求导时不能再将它视为候选参数 $\eta$ 的函数。

为简化推导，记当前轮的期望计数为
\begin{equation}
\begin{aligned}
a_{ik}^{(t)}&=\sum_{j=1}^N n_{ij}r_{ijk}^{(t)},\qquad
b_{kj}^{(t)}=\sum_{i=1}^M n_{ij}r_{ijk}^{(t)},\\
c_k^{(t)}&=\sum_{i=1}^M a_{ik}^{(t)}
=\sum_{j=1}^N b_{kj}^{(t)}.
\end{aligned}
\label{eq:lh-plsa-count-statistics}
\end{equation}
$a_{ik}^{(t)}$ 是话题 $k$ 产生单词 $i$ 的期望次数，$b_{kj}^{(t)}$ 是文本 $j$ 使用话题 $k$ 的期望次数，$c_k^{(t)}$ 是整个语料中分配给话题 $k$ 的期望总次数。于是
\[
Q=\sum_{i,k}a_{ik}^{(t)}\log\phi_{ik}
+\sum_{k,j}b_{kj}^{(t)}\log\theta_{kj}.
\]
引入列归一化约束的乘子 $\lambda_k$ 与 $\mu_j$，拉格朗日函数为
\[
\mathcal A=Q+
\sum_k\lambda_k\left(1-\sum_i\phi_{ik}\right)
+\sum_j\mu_j\left(1-\sum_k\theta_{kj}\right).
\]
若期望计数均为正，内点一阶条件给出
\[
\frac{a_{ik}^{(t)}}{\phi_{ik}}=\lambda_k,\qquad
\frac{b_{kj}^{(t)}}{\theta_{kj}}=\mu_j.
\]
将两式分别乘以 $\phi_{ik}$、$\theta_{kj}$ 并沿相应列求和，得到 $\lambda_k=c_k^{(t)}$ 与 $\mu_j=L_j$。因此更新式为
\begin{align}
\phi_{ik}^{(t+1)}
&=\frac{\sum_jn_{ij}r_{ijk}^{(t)}}
{\sum_{m=1}^M\sum_jn_{mj}r_{mjk}^{(t)}}
=\frac{a_{ik}^{(t)}}{c_k^{(t)}},
\label{eq:lh-plsa-m-step-phi}\\
\theta_{kj}^{(t+1)}
&=\frac{\sum_in_{ij}r_{ijk}^{(t)}}{L_j}
=\frac{b_{kj}^{(t)}}{L_j}.
\label{eq:lh-plsa-m-step-theta}
\end{align}
其中 $\sum_kb_{kj}^{(t)}=\sum_in_{ij}\sum_kr_{ijk}^{(t)}=L_j$，解释了第二个分母为何就是文本长度。$Q$ 是各概率列上凹函数的和，故上述解实现全局的 M 步最大化。若某些期望计数为零而对应总计数仍为正，归一化计数公式依然成立，零计数位置取零即可；这可由单列的 KL 散度非负性验证，无须在零点使用对数的内点求导。若 $c_k^{(t)}=0$，该轮 $Q$ 与 $\phi_{\cdot k}$ 无关，可保留原列，不能计算 $0/0$。

\begin{proposition}[EM 的似然单调性]
\label{prop:lh-plsa-em-monotonicity}
若当前观测对的模型概率为正，E 步精确计算式~\eqref{eq:lh-plsa-e-step}，M 步在约束下极大化式~\eqref{eq:lh-plsa-q-function}，则 $\ell(\eta^{(t+1)})\geq\ell(\eta^{(t)})$。由有限初值出发，对数似然值序列有有限极限，但这一结论不保证参数序列收敛，也不保证达到全局最大值。
\end{proposition}
\begin{proof}
对每个 $n_{ij}>0$ 的位置，任取话题分布 $q_{ij}=(q_{ijk})_{k=1}^K$，由对数函数的凹性，
\[
\log\sum_k\phi_{ik}\theta_{kj}
\geq\sum_kq_{ijk}\log\frac{\phi_{ik}\theta_{kj}}{q_{ijk}}.
\]
先在正概率情况下使用上式，零项再按极限约定处理。将右端乘以 $n_{ij}$ 并求和，记所得下界为 $F(\eta,q)$。当 $q_{ijk}=r_{ijk}^{(t)}$ 时，各正权重项中的比值在 $\eta=\eta^{(t)}$ 处相同，故下界取等号；固定 $q=r^{(t)}$ 后，$F$ 与 $Q$ 只差一个与 $\eta$ 无关的熵项。因此
\[
\ell(\eta^{(t+1)})\geq F(\eta^{(t+1)},r^{(t)})
\geq F(\eta^{(t)},r^{(t)})=\ell(\eta^{(t)}).
\]
又因每个 $p_{i\mid j}\leq1$，式~\eqref{eq:lh-plsa-loglikelihood} 中的条件对数似然不超过 $0$，故由单调有界性得到似然值的收敛。该论证只约束目标值，既没有证明参数的唯一性，也没有证明整个优化问题的凹性。
\end{proof}

分别固定 $\Phi$ 或 $\Theta$ 时，$\ell$ 是另一组参数的凹函数，但这不推出其关于两组参数联合凹。若 EM 参数序列收敛到所有参数均为正的内点，由更新映射的连续性，极限是 EM 不动点；将更新式代入似然的一阶导数可验证其满足归一化约束下的驻点条件。边界解需要单独检查相应的约束最优性条件，不能由似然值收敛直接断言是局部最大值。

\subsection{算法实施与手算例题}
\label{subsec:lh-plsa-em-implementation}

算法输入为非负整数矩阵 $X$、话题数 $K$、容差 $\varepsilon>0$ 和最大迭代次数 $T_{\max}$；输出为话题词分布 $\Phi$、文本话题分布 $\Theta$，以及需要联合模型时的 $\widehat\pi_j=L_j/n$。在移除空文本与全零词行后，可按以下步骤计算。
\begin{enumerate}
\item 为 $\Phi^{(0)}$ 和 $\Theta^{(0)}$ 选取严格正的初值，并分别按列归一化，计算初始条件对数似然。可以先随机产生正数再归一化；随机性来自初始化，随后的标准 EM 更新本身是确定的。
\item 对每个 $n_{ij}>0$ 的位置执行 E 步，累计式~\eqref{eq:lh-plsa-count-statistics} 的期望计数，再用式~\eqref{eq:lh-plsa-m-step-phi} 与式~\eqref{eq:lh-plsa-m-step-theta} 同时更新全部参数。当前轮的全部责任度均应使用旧参数，不能在累计过程中混用部分更新后的列。
\item 重新计算 $\ell^{(t+1)}=\ell(\Phi^{(t+1)},\Theta^{(t+1)})$。若增量非负且满足
\[
\frac{\ell^{(t+1)}-\ell^{(t)}}{\max\{1,|\ell^{(t)}|\}}
\leq\varepsilon,
\]
则返回当前参数；否则继续迭代。达到 $T_{\max}$ 时也返回当前近似解，并注明是否满足停止条件。若出现超出舍入误差的似然下降，应检查实现，不能将其当作正常收敛。
\end{enumerate}
停止条件检测的是似然增量，不能作为全局最优的证书。对不同初始化重复运行并比较最终似然，可以减少选到较差解的机会，但仍不能提供全局最优保证。若全部话题以完全相同的词分布和对称权重初始化，更新可能始终保持这种对称性，因此应避免用这样的初值期待算法自动区分话题。

令 $S=\#\{(i,j):n_{ij}>0\}$ 为非零计数位置的个数。一次 E 步及计数累计需要 $O(SK)$ 运算，归一化需要 $O(K(M+N))$；稠密情况下 $S=MN$。责任度可以逐个非零位置计算并立即累计，无须保存 $M\times N\times K$ 的数组，除存储输入外工作空间可为 $O(K(M+N))$。极小概率的计算可使用 $\log\phi_{ik}+\log\theta_{kj}$ 和 log-sum-exp 技巧，并检查每列的非负性与归一化。直接在更新后加任意常数并重新归一化会改变标准 EM 的 M 步，不能未经分析就沿用原来的单调性结论。

% 自拟例题：用精确分数演示一次 EM 更新，不是教材习题 18.3 的数值结果。
\begin{example}
\label{ex:lh-plsa-one-em-step}
设有两个单词、两个文本和两个话题，频数及初值为
\[
X=\begin{pmatrix}2&1\\1&2\end{pmatrix},\qquad
\Phi^{(0)}=\begin{pmatrix}3/4&1/4\\1/4&3/4\end{pmatrix},\qquad
\Theta^{(0)}=\begin{pmatrix}3/5&2/5\\2/5&3/5\end{pmatrix}.
\]
计算一轮 EM 更新，并检查词分布与文本话题分布的归一化。
\end{example}
\begin{solve}
两个文本长度都为 $3$，故 $\widehat\pi=(1/2,1/2)$。初始条件概率矩阵为
\[
B^{(0)}=\Phi^{(0)}\Theta^{(0)}
=\begin{pmatrix}11/20&9/20\\9/20&11/20\end{pmatrix}.
\]
只列出话题 $1$ 的责任度，话题 $2$ 的责任度逐项取补数：
\[
R_1^{(0)}=(r_{ij1}^{(0)})_{i,j}
=\begin{pmatrix}9/11&2/3\\1/3&2/11\end{pmatrix}.
\]
例如，$r_{111}^{(0)}=((3/4)(3/5))/(11/20)=9/11$，大于文本 $1$ 的先验话题比例 $3/5$，因为单词 $1$ 在话题 $1$ 中也更常见。将责任度分别乘以频数，话题 $1$ 的期望计数矩阵为
\[
\begin{pmatrix}18/11&2/3\\1/3&4/11\end{pmatrix}.
\]
它的行和为 $76/33$、$23/33$，列和为 $65/33$、$34/33$，总和为 $3$。归一化后得到
\[
\Phi^{(1)}=\begin{pmatrix}76/99&23/99\\23/99&76/99\end{pmatrix},\qquad
\Theta^{(1)}=\begin{pmatrix}65/99&34/99\\34/99&65/99\end{pmatrix}.
\]
两矩阵各列之和都为 $1$，而计数中的分数正是条件期望，并不要求取整数。新的 $B^{(1)}$ 对角元为 $5722/9801$，非对角元为 $4079/9801$，故条件对数似然由
\[
4\log(11/20)+2\log(9/20)\approx-3.98836
\]
增加至
\[
4\log(5722/9801)+2\log(4079/9801)\approx-3.90593.
\]
这验证了本轮更新的似然上升；它只是一次迭代的结果，尚不能视为收敛解。
\end{solve}

\begin{notealert}[type=caution,title={文本指标与新文本}]
基本 PLSA 为训练集合中的每个文本单独保存一列 $\theta_{\cdot j}$，并未给新文本的话题比例规定先验分布，也未直接给出生成未见文本指标的机制。若要表示一个新文本，可以固定学到的 $\Phi$，仅对新文本的混合比例 $\theta_{\cdot\mathrm{new}}$ 最大化其词项条件似然；责任度仍用式~\eqref{eq:lh-plsa-e-step}，再以新文本词频和长度更新该列。这是附加的参数拟合步骤，词表外单词也需预先处理，不能认为原模型已自动解决新文本推断。
\end{notealert}

\subsection{共现模型的 EM 算法}
\label{subsec:lh-plsa-cooccurrence-em}

教材还要求推导共现模型的 EM 算法。其联合对数似然为
\[
\ell_{\mathrm{co}}(\rho,\Phi,\Psi)
=\sum_{i,j}n_{ij}\log\left(\sum_k\rho_k\phi_{ik}\psi_{jk}\right).
\]
E 步由贝叶斯公式得到
\begin{equation}
\widetilde r_{ijk}^{(t)}=
\frac{\rho_k^{(t)}\phi_{ik}^{(t)}\psi_{jk}^{(t)}}
{\sum_{h=1}^K\rho_h^{(t)}\phi_{ih}^{(t)}\psi_{jh}^{(t)}}.
\label{eq:lh-plsa-cooccurrence-e-step}
\end{equation}
记 $\widetilde a_{ik}^{(t)}=\sum_jn_{ij}\widetilde r_{ijk}^{(t)}$、$\widetilde b_{kj}^{(t)}=\sum_in_{ij}\widetilde r_{ijk}^{(t)}$、$\widetilde c_k^{(t)}=\sum_{i,j}n_{ij}\widetilde r_{ijk}^{(t)}$。完全数据对数似然的条件期望为
\[
Q_{\mathrm{co}}=\sum_k\widetilde c_k^{(t)}\log\rho_k
+\sum_{i,k}\widetilde a_{ik}^{(t)}\log\phi_{ik}
+\sum_{k,j}\widetilde b_{kj}^{(t)}\log\psi_{jk}.
\]
三组参数分别在各自的概率单纯形上优化。由于 $\sum_k\widetilde c_k^{(t)}=n$，对三组归一化约束分别使用前述乘子法，得到 M 步
\begin{equation}
\rho_k^{(t+1)}=\frac{\widetilde c_k^{(t)}}{n},\qquad
\phi_{ik}^{(t+1)}=\frac{\widetilde a_{ik}^{(t)}}{\widetilde c_k^{(t)}},\qquad
\psi_{jk}^{(t+1)}=\frac{\widetilde b_{kj}^{(t)}}{\widetilde c_k^{(t)}}.
\label{eq:lh-plsa-cooccurrence-m-step}
\end{equation}
当 $\widetilde c_k^{(t)}=0$ 时，更新 $\rho_k^{(t+1)}=0$，而 $\phi_{\cdot k}$ 与 $\psi_{\cdot k}$ 可保留原来的归一化分布。其余停止条件、零概率处理与似然单调性的论证与生成模型相同。

若两种模型以命题~\ref{prop:lh-plsa-equivalence} 中相互对应的正参数初始化，式~\eqref{eq:lh-plsa-cooccurrence-e-step} 与式~\eqref{eq:lh-plsa-e-step} 给出相同责任度。共现模型完成一轮 M 步后，
\[
\pi_j^{(t+1)}=\sum_k\rho_k^{(t+1)}\psi_{jk}^{(t+1)}
=\frac{\sum_k\widetilde b_{kj}^{(t)}}{n}=\frac{L_j}{n},
\]
再由式~\eqref{eq:lh-plsa-cooccurrence-to-generative} 转换得到
\[
\theta_{kj}^{(t+1)}=
\frac{\rho_k^{(t+1)}\psi_{jk}^{(t+1)}}{\pi_j^{(t+1)}}
=\frac{\widetilde b_{kj}^{(t)}}{L_j}.
\]
这正是生成模型的更新式，说明两套算法的不同形式来自参数化的不同，而不是在学习两种不相容的概率规律。实际使用时，$\Phi$ 的高概率词可用于解释话题，$\Theta$ 可用于表示训练文本的混合组成；具体话题名称及话题数的选择仍需结合任务与拟合结果判断。


% 教材来源：机器学习/textbooks/机器学习（分章原文）/机器学习-20.pdf。
% 对应教材第 20 章，原书第 385--414 页；保留本文的自动章节编号。
% 变分推导补充核对：Blei, Ng, Jordan, Latent Dirichlet Allocation,
% JMLR 3 (2003), 993--1022, §§5.1、5.3--5.4 及附录 A。
% https://www.jmlr.org/papers/volume3/blei03a/blei03a.pdf




